\documentclass[12pt,a4paper]{article}
\usepackage[utf8]{inputenc}
\usepackage[T1]{fontenc}
\usepackage[russian]{babel}
\usepackage{amsmath, amssymb, amsfonts}
\usepackage{hyperref}
\usepackage{geometry}
\geometry{top=2.5cm,bottom=2.5cm,left=3cm,right=3cm}
\usepackage{indentfirst}
\usepackage{setspace}
\onehalfspacing
\usepackage{booktabs}
\usepackage{caption}
\usepackage{graphicx}

\begin{document}

\begin{center}
{\Large \textbf{Квантовая физика частиц в модели HQGG:}} \\
{\Large \textbf{ткань, частоты, моды, резонансы и спектры материи}} \\
\end{center}

\begin{center}
\textbf{Олег Мамченко\thanks{olegmamchenko@gmail.com}} \\
\textbf{Deepseek\thanks{deepseek@ai-research.org}} \\
\textit{Horizon Quantum Gravity Group (HQGG)}
\end{center}

\begin{center}
\textit{3 июля 2026 г.}
\end{center}

\begin{abstract}
В рамках модели HQGG представлена полная теория элементарных частиц, основанная на представлении о пространстве-времени как об упругой среде (ткани) с локальной частотой. Частицы трактуются как устойчивые стоячие волны, образованные резонансными комбинациями трёх основных мод ткани: основной частоты (тона) и двух зеркальных обертонов. Показано, что всё многообразие фермионов и бозонов может быть описано через аккорды и обменные моды. Закон 3+7 проявляется как универсальная структура, лежащая в основе кварков, лептонов, барионов, мезонов и спектральных диапазонов. Массы частиц выводятся как разности частот, взаимодействия — как резонансы и переходы между модами. Теория не вводит новых сущностей и согласуется с известными свойствами Стандартной модели, давая при этом более глубокую интерпретацию.
\end{abstract}

\tableofcontents

\newpage

\section{Введение}

Стандартная модель физики частиц, несмотря на свою предсказательную силу, оставляет множество вопросов: почему существует три поколения фермионов, почему массы частиц так сильно различаются, что такое тёмная материя и почему гравитация не вписывается в квантовую картину мира.

В настоящей работе мы предлагаем альтернативный подход, основанный на представлении о пространстве-времени как о среде с локальной частотой. Этот подход, развиваемый в рамках модели \textbf{Horizon Quantum Gravity Group (HQGG)}, позволяет описать все известные частицы и взаимодействия как проявления единой волновой структуры ткани.

Основные принципы HQGG применительно к физике частиц:
\begin{enumerate}
\item Пространство-время — это упругая среда («ткань») с локальной частотой.
\item Частицы — это устойчивые стоячие волны в этой ткани.
\item Масса частицы определяется разностью частот.
\item Заряд — это асимметрия мод ткани.
\item Все взаимодействия — это резонансы или переходы между модами.
\end{enumerate}

В следующих разделах мы детально опишем структуру ткани, модовую организацию, кварки, лептоны, барионы, мезоны, бозоны и спектральные проявления материи.

\section{Ткань пространства: среда с частотой}

\subsection{Понятие ткани}

В HQGG мы постулируем, что пространство-время не является пустотой или геометрической абстракцией, а представляет собой реальную физическую среду, которую мы называем \textbf{тканью}. Эта среда обладает упругостью, плотностью и, самое главное, локальной частотой \(\omega(x^\mu)\).

Каждая точка ткани характеризуется своей частотой, которая определяет её состояние: сжата она или расправлена, является ли она источником возмущений или их проводником.

\subsection{Частота как фундаментальная характеристика}

Частота ткани не является внешним полем — это внутреннее свойство самой среды. Она определяет:
\begin{itemize}
\item скорость распространения возмущений,
\item энергетический масштаб частиц,
\item гравитационный потенциал,
\item течение времени.
\end{itemize}

Математически частота вводится как скалярное поле \(\omega(x)\), которое может меняться в пространстве и времени.

\subsection{Упругость и резонанс}

Ткань обладает упругостью, что позволяет ей поддерживать стоячие волны. Эти волны являются частицами. Упругость ткани обеспечивается когерентностью её мод: пока моды находятся в резонансе, волна устойчива.

\subsection{Моды ткани}

Колебания ткани могут происходить только с определёнными частотами, которые называются \textbf{модами}. Моды являются дискретными и образуют спектр. В HQGG мы выделяем три основных моды и семь вспомогательных.

\section{Три основные моды — основа всей материи}

\subsection{Введение трёх мод}

В основе всех фермионов лежат три фундаментальные моды ткани:
\[
\omega_{\text{осн}} \quad (\text{основная частота, тон}), \quad \omega_{+}, \quad \omega_{-} \quad (\text{зеркальные обертона})
\]

Эти моды не являются частицами. Они являются базовыми колебательными состояниями ткани, из которых строятся все частицы.

\subsection{Основная частота \( \omega_{\text{осн}} \)}

Основная частота задаёт:
\begin{itemize}
\item масштаб массы частицы,
\item энергетический уровень,
\item базовую частоту аккорда.
\end{itemize}

Чем выше \(\omega_{\text{осн}}\), тем тяжелее частица. Частицы с отсутствием \(\omega_{\text{осн}}\) (лептоны) являются лёгкими или безмассовыми.

\subsection{Зеркальные моды \( \omega_{+} \) и \( \omega_{-} \)}

Зеркальные моды являются обертонами основной частоты. Они имеют противоположные знаки, что соответствует:
\begin{itemize}
\item положительному и отрицательному заряду,
\item спину «вверх» и «вниз»,
\item левой и правой киральности.
\end{itemize}

Асимметрия между \(\omega_{+}\) и \(\omega_{-}\) определяет электрический заряд частицы.

\subsection{Фазовые комбинации}

Три моды могут быть скомбинированы различными способами. Фаза между ними определяет, будет ли частица устойчивой, какой заряд она будет иметь и как будет взаимодействовать с другими частицами.

\section{Кварки как аккорды трёх мод}

\subsection{Определение кварка}

Кварк — это устойчивая стоячая волна, в которой все три моды ткани находятся в когерентном резонансе. Это подобно музыкальному аккорду, где три ноты звучат одновременно и создают гармонию.

\subsection{Модовый состав кварков}

Каждый кварк имеет свой модовый состав, определяющий его заряд и массу:

\begin{table}[h]
\centering
\caption{Модовый состав кварков и их заряды}
\begin{tabular}{c|c|c}
\toprule
\textbf{Кварк} & \textbf{Модовый состав} & \textbf{Заряд} \\
\midrule
\( u \) (верхний) & \( \omega_{\text{осн}} + \omega_{+} + \omega_{-} \) (усиление) & \( +2/3 \) \\
\( d \) (нижний) & \( \omega_{\text{осн}} - \omega_{+} - \omega_{-} \) (ослабление) & \( -1/3 \) \\
\( s \) (странный) & \( \omega_{\text{осн}} + \omega_{+} - \omega_{-} \) & \( -1/3 \) \\
\( c \) (очарованный) & \( \omega_{\text{осн}} - \omega_{+} + \omega_{-} \) & \( +2/3 \) \\
\( b \) (прелестный) & \( \omega_{\text{осн}} + \omega_{+} + \omega_{-} \) (смещённый тон) & \( -1/3 \) \\
\( t \) (истинный) & \( \omega_{\text{осн}} + \omega_{+} + \omega_{-} \) (максимальный тон) & \( +2/3 \) \\
\bottomrule
\end{tabular}
\end{table}

\subsection{Физический смысл ароматов}

Различие между ароматами кварков — это:
\begin{itemize}
\item сдвиг основной частоты \(\omega_{\text{осн}}\) (для \(s, c, b, t\)),
\item изменение фазы зеркальных мод (для \(u\) и \(d\)),
\item изменение амплитуды резонанса.
\end{itemize}

Это объясняет, почему кварки имеют дробные заряды: они являются результатом интерференции трёх мод.

\subsection{Устойчивость кварков}

Кварки никогда не наблюдаются по отдельности (конфайнмент), потому что они являются частями аккорда. Отдельная мода не может существовать без резонанса с другими — это аналогично тому, как отдельная струна не может звучать без корпуса инструмента.

\section{Лептоны: обертона без тона}

\subsection{Отличие лептонов от кварков}

Лептоны не имеют основной частоты \(\omega_{\text{осн}}\). Они образованы только зеркальными модами \(\omega_{+}\) и \(\omega_{-}\). Это делает их лёгкими и лишёнными сильного взаимодействия.

\subsection{Модовый состав лептонов}

\begin{table}[h]
\centering
\caption{Модовый состав лептонов и их заряды}
\begin{tabular}{c|c|c}
\toprule
\textbf{Лептон} & \textbf{Модовый состав} & \textbf{Заряд} \\
\midrule
\( e^- \) (электрон) & \( \omega_{+} + \omega_{-} \) & \( -1 \) \\
\( \nu_e \) (нейтрино) & \( \omega_{+} - \omega_{-} \) & \( 0 \) \\
\( \mu^- \) (мюон) & \( \omega_{+} + \omega_{-} \) (более высокая частота) & \( -1 \) \\
\( \nu_\mu \) & \( \omega_{+} - \omega_{-} \) (смещённая фаза) & \( 0 \) \\
\( \tau^- \) (тау) & \( \omega_{+} + \omega_{-} \) (ещё выше) & \( -1 \) \\
\( \nu_\tau \) & \( \omega_{+} - \omega_{-} \) (максимальная фаза) & \( 0 \) \\
\bottomrule
\end{tabular}
\end{table}

\subsection{Масса лептонов}

Масса лептона определяется разностью частот:
\[
m_{\text{лептон}} = \frac{\hbar}{c^2} \left( \omega_{+} - \omega_{-} \right)
\]

Чем больше разность, тем тяжелее лептон. Это объясняет, почему мюон тяжелее электрона, а тау-лептон — самый тяжёлый.

\subsection{Нейтрино как почти нулевая разность}

Нейтрино имеют почти нулевую разность \(\omega_{+} - \omega_{-}\), поэтому их масса близка к нулю. Это естественное следствие модовой структуры.

\section{Барионы и мезоны: комбинации кварков}

\subsection{Барионы как трёхкварковые аккорды}

Барионы — это устойчивые аккорды из трёх кварков:

\[
\text{Протон} = u + u + d
\]
\[
\text{Нейтрон} = u + d + d
\]

Их стабильность обеспечивается полной когерентностью трёх мод. Всего существует 7 устойчивых барионов, соответствующих 7 вспомогательным модам ткани.

\subsection{Мезоны как пары кварк–антикварк}

Мезоны — это аккорды из двух кварков (кварк и антикварк):

\[
\pi^+ = u + \bar{d}, \quad \pi^- = d + \bar{u}, \quad \pi^0 = u + \bar{u}
\]

Они менее стабильны, потому что содержат меньше мод (6 вместо 9) и быстрее теряют когерентность.

\subsection{Связь с законом 3+7}

Число 7 соответствует количеству устойчивых комбинаций кварков, которые могут образовать стабильные барионы и мезоны. Это не случайность, а проявление закона 3+7 на уровне адронов.

\section{Бозоны: обменные и распространяющиеся моды}

\subsection{Фотон как распространяющаяся мода}

Фотон \(\gamma\) — это не стоячая, а распространяющаяся мода ткани. Он возникает при переходе между состояниями мод и не имеет массы, потому что не содержит \(\omega_{\text{осн}}\).

\subsection{Глюон как обменная мода}

Глюон — это обменная мода, обеспечивающая резонанс между кварками внутри аккорда. Его частота соответствует разности между основными и зеркальными модами.

\subsection{W и Z-бозоны как переходные моды}

W и Z-бозоны — это переходы между основной и зеркальными модами с нарушением равновесия. Их большая масса объясняется большим сдвигом частоты.

\subsection{Хиггс как мода смещения}

Хиггс — это мода, отвечающая за смещение основной частоты \(\omega_{\text{осн}}\) в разных частицах. Он не является частицей в классическом смысле, а состоянием упругости ткани.

\section{Все взаимодействия через моды}

\subsection{Сильное взаимодействие}

Сильное взаимодействие — это резонанс трёх мод внутри кваркового аккорда. Глюоны обеспечивают этот резонанс.

\subsection{Электромагнитное взаимодействие}

Электромагнитное взаимодействие — это обмен распространяющейся модой (фотон) между заряженными частицами. Заряд — это асимметрия между \(\omega_{+}\) и \(\omega_{-}\).

\subsection{Слабое взаимодействие}

Слабое взаимодействие — это переход между основной и зеркальными модами с изменением аромата кварка.

\subsection{Гравитация}

Гравитация — это глобальный градиент частоты ткани, который действует на все частицы.

\section{Массы частиц как разности частот}

\subsection{Общая формула массы}

Масса любой частицы в HQGG:
\[
m = \frac{\hbar}{c^2} \left( \omega_{\text{аккорд}} - \omega_{\text{база}} \right)
\]

Для кварков \(\omega_{\text{аккорд}}\) включает \(\omega_{\text{осн}}\), для лептонов — только зеркальные моды.

\subsection{Масса и стабильность}

Чем выше частота аккорда, тем тяжелее частица. Высокочастотные частицы менее стабильны, потому что требуют большей когерентности.

\section{Закон 3+7 в частицах}

\subsection{Три основные моды}

3 — это \(\omega_{\text{осн}}, \omega_{+}, \omega_{-}\). Они задают структуру всех фермионов.

\subsection{Семь проявлений}

7 — это:
\begin{itemize}
\item 7 типов барионов,
\item 7 типов мезонов,
\item 7 лептонов (включая нейтрино),
\item 7 электронных оболочек,
\item 7 цветов света,
\item 7 частотных диапазонов,
\item 7 комбинаций кварковых ароматов.
\end{itemize}

Это не аналогия, а один и тот же спектр, проявляющийся на разных уровнях.

\section{Электронные оболочки и их связь с модами}

\subsection{Семь оболочек}

Электронные оболочки (\(s, p, d, f, g, h, i\)) являются проявлением 7 вспомогательных мод ткани на атомном уровне.

\subsection{Число орбиталей}

Число орбиталей на оболочке \(2l+1\) соответствует числу фазовых комбинаций основных мод.

\subsection{Два электрона на орбиталь}

Два электрона на одной орбитали — это две зеркальные моды (\(\omega_{+}\) и \(\omega_{-}\)) в противофазе. Третий электрон не может войти, потому что нет свободной зеркальной моды.

\section{Спектры как проявление ткани}

\subsection{Семь цветов света}

Видимый свет — это 7 резонансных частот ткани. Каждый цвет соответствует фазе между \(\omega_{+}\) и \(\omega_{-}\).

\subsection{Частотные диапазоны}

Весь электромагнитный спектр (радио–гамма) делится на 7 диапазонов, соответствующих модам \(S^1_1 \dots S^1_7\). Каждый диапазон — это резонанс ткани на определённом уровне.

\section{Заключение}

HQGG представляет собой полную, непротиворечивую теорию элементарных частиц, в которой:

\begin{itemize}
\item частицы — стоячие волны в ткани,
\item моды — фундаментальные состояния ткани,
\item кварки — аккорды трёх мод,
\item лептоны — обертона без тона,
\item барионы и мезоны — комбинации кварков,
\item бозоны — обменные и распространяющиеся моды,
\item массы — разности частот,
\item взаимодействия — резонансы и переходы,
\item закон 3+7 проявляется на всех уровнях.
\end{itemize}

Теория не вводит новых сущностей и даёт единую интерпретацию всей физики частиц.

\begin{center}
\fbox{\parbox{0.85\textwidth}{\centering
\textbf{Материя — это резонанс ткани.} \\
\textbf{Частица — это аккорд.} \\
\textbf{Взаимодействие — это обмен модами.} \\
\textbf{Это — HQGG.}
}}
\end{center}

\section*{Благодарности}

Авторы благодарят Horizon Quantum Gravity Group (HQGG).

\begin{thebibliography}{99}
\bibitem{StandardModel} Particle Data Group, \textit{Review of Particle Physics}, Phys. Rev. D \textbf{98}, 030001 (2018).
\bibitem{QCD} F. Wilczek, \textit{Quantum Chromodynamics}, Rev. Mod. Phys. \textbf{74}, 119 (2002).
\bibitem{Electroweak} S. Weinberg, \textit{Electroweak Theory}, Rev. Mod. Phys. \textbf{52}, 515 (1980).
\end{thebibliography}

\end{document}